%%=============================================================================
%% Methodologie
%%=============================================================================

\chapter{\IfLanguageName{dutch}{Methodologie}{Methodology}}%
\label{ch:methodologie}

Dit hoofdstuk beschrijft het plan van aanpak waarmee de onderzoeksvragen uit de
inleiding worden beantwoord. Het bouwt rechtstreeks voort op de bevindingen van de
literatuurstudie in Hoofdstuk~\ref{ch:stand-van-zaken} en vertaalt die naar een
concreet onderzoeksontwerp. Het onderzoek wordt opgevat als een vergelijkende
studie: er bestaat geen op voorhand vaststaand ``beste'' model voor onze casus,
dus worden meerdere modellen en twee taakvormen onder gelijke omstandigheden tegen
elkaar afgewogen. Dit hoofdstuk vormt een eerste versie van die methodologie en
schetst vooral het geplande onderzoeksproces; de uiteindelijke resultaten en de
conclusie volgen in de latere hoofdstukken. Omdat het onderzoek een exploratief en
iteratief karakter heeft, kan de hier beschreven aanpak tijdens de uitvoering nog
verfijnd en uitgebreid worden wanneer praktische keuzes of tussentijdse experimenten
daarom vragen.

\section{Onderzoeksopzet}

Om de centrale onderzoeksvraag te beantwoorden, wordt het werk opgedeeld in een
aantal opeenvolgende fasen. Elke fase heeft een eigen doelstelling en levert een
concrete tussenstap (deliverable) op die als input dient voor de volgende fase:

\begin{enumerate}
    \item \textbf{Literatuurstudie en afbakening} --- de theoretische basis en de
    keuze van de onderzoeksrichting.
    \item \textbf{Dataverzameling en opbouw van de dataset} --- het aanleggen van een
    eigen beelddataset bij Delhaize Ledeberg.
    \item \textbf{Voorverwerking en datasetsplitsing} --- het klaarmaken van de data
    voor training en evaluatie.
    \item \textbf{Selectie en training van de modellen} --- het trainen van een
    shortlist aan classificatiemodellen en een detectiemodel via transfer
    learning.
    \item \textbf{Experimentele evaluatie} --- het meten en vergelijken van de
    prestaties onder verschillende omstandigheden.
    \item \textbf{Analyse en Proof of Concept} --- het selecteren van de beste aanpak
    en het bouwen van een werkende proof of concept.
\end{enumerate}

Een rode draad doorheen deze fasen is het tweeledige doel uit de probleemstelling:
enerzijds het grofmazige minimumdoel (de juiste broodzakcategorie bepalen) en
anderzijds het fijnmazigere onderzoeksdoel (de specifieke broodsoort onderscheiden).
Beide niveaus worden met dezelfde dataset en grotendeels dezelfde modellen
onderzocht.

\section{Literatuurstudie en afbakening}

De eerste fase bestond uit de literatuurstudie, die in Hoofdstuk~\ref{ch:stand-van-zaken}
is uitgewerkt en hier niet wordt herhaald. Deze fase had als doel de bestaande
technieken voor brood- en voedselherkenning in kaart te brengen en de voornaamste
methodologische keuzes te onderbouwen. De belangrijkste conclusies sturen de rest
van dit hoofdstuk rechtstreeks aan: beeldclassificatie is theoretisch de meest
directe taak voor onze opstelling, maar objectdetectie wordt parallel getest;
transfer learning is aangewezen bij een beperkte eigen dataset; en er bestaat geen
publieke dataset die de broodsoorten van Delhaize Ledeberg als afzonderlijke klassen
bevat. Die laatste vaststelling verantwoordt meteen de volgende fase, waarin een
eigen dataset wordt opgebouwd.

\section{Dataverzameling en opbouw van de dataset}

In deze fase wordt een eigen beelddataset aangelegd met foto's van de broodsoorten
uit het assortiment van Delhaize Ledeberg. De opnames gebeuren ter plaatse, in een
opstelling die de uiteindelijke toepassing aan de broodsnijmachine zo goed mogelijk
benadert. Het exacte aantal broodsoorten dat in de dataset terechtkomt, kan pas op
de dag van de opnames definitief worden vastgesteld, aangezien dit afhangt van het
assortiment en de beschikbaarheid in de winkel op dat moment. Ook de beschikbare tijd
op de opnamedag bepaalt mee hoeveel data effectief verzameld kan worden.

Als streefdoel wordt gemikt op ongeveer 75 tot 100 originele foto's per broodsoort.
Per individueel fysiek brood worden ongeveer vijf foto's genomen, wat betekent dat er
per broodsoort zo'n vijftien tot twintig verschillende exemplaren nodig zijn. Deze
keuze volgt rechtstreeks uit de literatuurstudie, waar vooral de variatie tussen
verschillende fysieke exemplaren van dezelfde soort belangrijk bleek.

Om de probleemdomein-deelvragen over licht en bedekking te kunnen beantwoorden, wordt
een bewust onderscheid gemaakt tussen de manier waarop die variatie in de training en
in de evaluatie terechtkomt. De hoofddataset bestaat hoofdzakelijk uit normale,
realistische foto's van broden aan de broodsnijmachine. Tijdens het trainen wordt
variatie in onder meer belichting, kleur en gedeeltelijke bedekking kunstmatig
nagebootst met data-augmentatie, wat het model robuuster maakt zonder dat daarvoor
grote aantallen extra opnames nodig zijn~\autocite{Shorten2019}.

Synthetische augmentatie alleen volstaat echter niet om die deelvragen geloofwaardig
te beantwoorden: een helderheidsaanpassing of een kunstmatig gewist beelddeel
benadert de werkelijkheid slechts ruw, terwijl echt winkellicht ook schaduwen en
reflecties verandert~\autocite{Olakanmi2023} en een echte hand of zak het brood op
een natuurlijke manier bedekt~\autocite{Ciapas2023}. Daarom worden tijdens de
dataverzameling bijkomend een beperkt aantal echte foto's genomen onder afwijkende
lichtomstandigheden en met echte gedeeltelijke bedekking. Deze beelden worden niet aan
de trainingsdata toegevoegd, maar apart gehouden als \emph{challenge-testset} waarmee
de robuustheid van de modellen onder realistische omstandigheden wordt geëvalueerd.

Elke specifieke broodsoort wordt gelabeld op soortniveau en tegelijk gekoppeld aan de
correcte broodzakcategorie. Daardoor kan dezelfde dataset zowel op het fijnmazige
niveau (broodsoort) als op het grofmazige niveau (broodzakcategorie) worden
geëvalueerd, zonder dat er opnieuw data verzameld moet worden. Voor de objectdetectie
worden op dezelfde ruwe beelden bijkomend bounding boxes rond de broden geannoteerd;
de onderliggende foto's blijven dus identiek, enkel de annotatie verschilt.

\section{Voorverwerking en datasetsplitsing}

Voor de verzamelde beelden bruikbaar zijn, worden ze eerst gecontroleerd op
kwaliteit (scherpte, belichting, correcte labels) en consistent voorverwerkt, onder
meer door ze te herschalen naar de invoerresolutie van de gekozen modellen. De normale
hoofddataset wordt vervolgens opgesplitst in een trainings-, validatie- en testset
volgens een 70/15/15-verdeling, waarbij 70 procent naar de trainingsset gaat en
telkens 15 procent naar de validatie- en testset.

Bij het splitsen wordt een belangrijke methodologische voorzorg genomen: de verdeling
gebeurt op het niveau van het fysieke brood en niet willekeurig per foto. Alle foto's
van eenzelfde exemplaar belanden dus volledig in dezelfde subset. Zo wordt vermeden
dat sterk gelijkende opnames van hetzelfde brood zowel in de trainings- als in de
testset voorkomen, wat de prestaties kunstmatig zou opdrijven (\emph{data leakage}).
Deze opzet garandeert dat het model op de testset wordt beoordeeld op broden die het
tijdens de training nooit heeft gezien.

Naast deze normale testset staat de eerder vermelde challenge-testset met echte licht-
en occlusiebeelden. Die blijft volledig buiten de trainings- en validatieset en dient
uitsluitend voor de eindevaluatie van de robuustheid. Waar praktisch mogelijk worden
hiervoor fysieke broodexemplaren gebruikt die niet in de trainings- of validatieset
voorkomen, zodat ook hier geen data leakage optreedt.

Data-augmentatie wordt uitsluitend op de trainingsset toegepast. De validatieset, de
normale testset en de challenge-testset bestaan dus steeds uit echte,
niet-geaugmenteerde beelden, zodat elke evaluatie op realistische beelden gebeurt. De
augmentatietechnieken sluiten aan bij de deelvragen rond licht en bedekking:
geometrische transformaties, variatie in helderheid en kleur, en het willekeurig
wissen van beelddelen~\autocite{Shorten2019}. Op die manier is augmentatie een
trainingsstrategie om robuustheid op te bouwen, terwijl de echte challengebeelden de
onafhankelijke toets van die robuustheid vormen.

\section{Selectie en training van de modellen}

In deze fase worden de modellen getraind en vergeleken. De gemeenschappelijke
basisstrategie is transfer learning: er wordt vertrokken van CNN's die al op een
grote dataset zoals ImageNet zijn voorgetraind, waarna ze op de eigen brooddataset
worden bijgesteld. Op die manier blijft ook een relatief kleine dataset werkbaar.
Waar praktisch haalbaar wordt het effect van domeinspecifieke fine-tuning onderzocht,
door niet enkel de classificatielaag maar ook diepere lagen bij te stellen.

Voor beeldclassificatie wordt gewerkt met een shortlist van architecturen in plaats
van op voorhand één model als beste aan te duiden. Die aanpak is methodologisch
onderbouwd: \textcite{AbouBaker2022} tonen dat voorgetrainde netwerken sterk kunnen
verschillen in de afweging tussen nauwkeurigheid, trainingstijd en modelgrootte,
waardoor het zinvol is meerdere kandidaten te vergelijken. Deze voorlopige shortlist
bevat zowel een compacte, snelle architectuur (MobileNet) als een krachtigere backbone
(DenseNet), aangevuld met VGG-16 als brood-specifieke referentie uit
\textcite{Tantiphanwadi2022}. De keuze voor DenseNet sluit aan bij \textcite{Boyd2024},
die onder zeven vergeleken CNN-architecturen DenseNet als sterkste model voor
fijnmazige voedselherkenning aanwijzen. De samenstelling van deze shortlist ligt niet
definitief vast en kan tijdens de uitvoering nog worden bijgesteld als praktische
bevindingen daarom vragen.

Voor objectdetectie wordt een model uit de YOLO-familie onderzocht, dat één brood in
het beeld lokaliseert en tegelijk classificeert. De precieze YOLO-variant wordt tijdens
de uitvoering gekozen op basis van beschikbaarheid, prestaties en praktische
inzetbaarheid. Om een eerlijke vergelijking mogelijk te maken, worden de classificatie-
en detectiemodellen zoveel mogelijk onder vergelijkbare voorwaarden en op dezelfde
onderliggende beelden getraind.

\section{Experimentele evaluatie}

Zodra de modellen getraind zijn, worden ze systematisch geëvalueerd en onderling
vergeleken. Voor de classificatiemodellen wordt niet enkel naar de accuracy gekeken,
maar ook naar precision, recall en de F1-score. Omdat het om een meerklassenprobleem
met veel gelijkende broodsoorten gaat, worden deze maten als macrogemiddelde berekend,
zodat elke klasse even zwaar meetelt ongeacht het aantal beelden. Een confusion matrix
maakt daarbij zichtbaar welke broodsoorten het vaakst met elkaar worden verward. Voor
het detectiemodel wordt de \emph{mean average precision} (mAP) als primaire maat
gebruikt, aangevuld met precision en recall; afhankelijk van wat de gekozen
YOLO-implementatie ondersteunt, wordt mAP gerapporteerd bij een IoU-drempel van 0,5
(mAP@0.5) en eventueel over het bereik 0,5--0,95 (mAP@0.5:0.95).

Naast deze kwaliteitsmaten worden ook de inferentietijd en de modelgrootte gemeten.
Deze praktische criteria zijn essentieel omdat het systeem aan de broodsnijmachine
snel moet reageren en de klant niet lang mag wachten; de uiteindelijke modelkeuze mag
dan ook niet blind op één metriek steunen, maar moet nauwkeurigheid afwegen tegen
snelheid en inzetbaarheid.

De evaluatie verloopt in stappen die elk een onderzoeksdeelvraag bedienen. Eerst
worden de prestaties op de normale testset bepaald; die vormen de referentie waartegen
alles wordt afgezet. Vervolgens worden dezelfde relevante metrics berekend op de echte
licht-challengeset en daarna op de echte occlusie-challengeset. Door de prestatiedaling
ten opzichte van de normale testset te vergelijken, kan de impact van variërende
belichting en van gedeeltelijke bedekking afzonderlijk worden gekwantificeerd. Deze
vergelijking wordt waar mogelijk uitgevoerd voor beide taakvormen (classificatie en
detectie) en op beide niveaus (broodzakcategorie en specifieke broodsoort), zodat ook
de bijdrage van transfer learning en fine-tuning zichtbaar wordt.

\section{Analyse en Proof of Concept}

In de laatste fase worden de resultaten van alle experimenten samengebracht en
geanalyseerd. Op basis van de vergelijking tussen de modellen en taakvormen wordt de
best presterende en praktisch meest haalbare aanpak geselecteerd. Daarmee wordt een
werkende proof of concept gebouwd die via een camera een brood analyseert en minstens
de juiste broodzakcategorie adviseert, met daarbovenop de vraag hoever de herkenning
tot op het niveau van de specifieke broodsoort kan gaan. Het hoofdstuk mondt uit in een
technisch onderbouwd advies over de praktische haalbaarheid van automatische
broodherkenning voor Delhaize Ledeberg, dat in de conclusie wordt teruggekoppeld aan
de onderzoeksvragen.